% !TEX root = ../main.tex
% Draft 10: claim/source map in notes/DRAFT10_REVISION.md. Performance claims remain pending.
\section{Introduction}
\label{sec:introduction}
% Native vector schematic; no synthetic or illustrative performance data.
\begin{minipage}{\textwidth}
\centering
\captionsetup{type=figure,hypcap=false}
\begin{tikzpicture}[x=1cm,y=1cm,>=Latex,font=\small]
\definecolor{targetgreen}{RGB}{35,135,107}
\definecolor{protectred}{RGB}{186,76,65}
\definecolor{policyblue}{RGB}{36,105,145}
\foreach \x/\title in {0/1. Identify the object,5.65/2. Recover its surface,11.3/3. Authorize contact}{
 \fill[black!3,rounded corners=3pt] (\x,0) rectangle +(5.4,3.1);
 \node[anchor=north,font=\small\bfseries] at (\x+2.7,2.98) {\title};
 \draw[black!25,line width=1pt] (\x+.35,.85)--(\x+5.05,.85);
 \fill[targetgreen!22] (\x+2.2,1.5) circle (.63);
 \draw[targetgreen,line width=1.2pt] (\x+2.2,1.5) circle (.63);
 \fill[protectred!55] (\x+3.2,.85) rectangle +(0.65,1.25);
}
\draw[policyblue,dashed,thick] (2.6,1.5) ellipse (1.65 and 1.0);
\node[policyblue,font=\footnotesize] at (1.25,2.3) {coarse target bound};
\node[align=center,font=\footnotesize] at (2.7,.37) {The right identity can still\\include the wrong obstacle.};
\foreach \a in {0,20,...,340}{
 \fill[targetgreen] ({7.85+.60*cos(\a)},{1.5+.60*sin(\a)}) circle (.037);
}
\draw[protectred,thick] (8.85,.85) rectangle +(0.65,1.25);
\node[targetgreen,font=\footnotesize] at (7.55,2.32) {observed target};
\node[align=center,font=\footnotesize] at (8.35,.37) {Attention seeds the object;\\depth preserves the boundary.};
\draw[policyblue,line width=3pt] (12.55,2.3)--(12.55,1.5)--(12.85,1.5);
\draw[policyblue,line width=3pt] (14.4,2.3)--(14.4,1.5)--(14.15,1.5);
\fill[targetgreen] (12.88,1.5) circle (.085);
\fill[targetgreen] (14.12,1.5) circle (.085);
\node[policyblue,font=\footnotesize] at (13.5,2.48) {designated fingers};
\node[align=center,font=\footnotesize] at (14,.37) {Permission is link- and state-specific.\\Other surfaces remain protected.};
\end{tikzpicture}
\caption{\textbf{Identity, extent, and permission are different decisions.}
\sys\ uses internal noun attention to seed an observed 3D surface graph. The selected surface defines the intended extent of a contact exception; a coarse bound does not. The final panel illustrates link-specific contact permission; after grasping, the target becomes attached geometry checked against the environment. All panels illustrate the proposed mechanism, not measured execution outcomes.}
\label{fig:teaser}
\end{minipage}


A robot reaching for an apple must distinguish two geometrically similar events: its fingers touching the apple and its
wrist striking a nearby object. The distinction depends on the instruction, the robot link, and the stage of manipulation.
After grasping, the apple itself becomes moving geometry that must clear the environment; after release, it becomes part
of the scene again. This changing role of the same visible surface is a central difficulty in connecting
vision-language-action (VLA) policies to geometric collision avoidance.

VLAs such as $\pi_0$, $\pi_{0.5}$, and GR00T N1 generate action chunks from visual observations and language
instructions~\cite{black2024pi0,intelligence2025pi05,nvidia2025gr00t}. Their task-conditioned behavior makes them a useful
source of manipulation proposals, but successful task completion alone does not establish collision-free
execution~\cite{fan2026safevlabench}. Depth-based distance fields provide an explicit account of surrounding surfaces
for motion generation~\cite{sundaralingam2023curobo,sundaralingam2026curobov2}. Geometry alone, however, does not specify
which of those surfaces a particular link may touch. Treating the target as an obstacle can prevent the grasp; removing
it from a shared collision field also removes it for links that should still avoid it. The missing connection is from
\emph{visual task relevance} to \emph{spatially and temporally selective contact constraints}.

Prior work provides several parts of this connection. VLSA uses external visual grounding to configure a safety
filter~\cite{hu2026vlsa}, and MoManipVLA excludes a segmented target from its distance-field reconstruction in its
real-world pipeline~\cite{wu2025momanipvla}. KNOWS uses VLA attention to select a target among objects supplied by a
segmentation model~\cite{park2026knows}. Thus, attention-based target selection and target exclusion are established
ideas. What remains to be addressed here is how to recover the target's spatial extent directly from VLA signals and
depth, and use that geometry to determine contact permissions for individual links as manipulation proceeds.

Our approach separates the roles of attention and geometry: \emph{attention selects a candidate object, while depth
defines its observed extent}. We present \sys, which grounds action-to-image attention in a multi-view point cloud
without an external detector or segmentation network. Attention seeds geometric clustering rather than masking the
collision map, so weakly attended surfaces remain available as obstacles. The selected target then defines a local
contact exception for the active hand. Other modeled links continue to query the scene field containing the target.
When a grasp is confirmed, the target is represented as geometry attached to the hand and checked against the
environment. This couples object grounding to the geometric consequences of interaction (Fig.~\ref{fig:overview}).

To control when target identities may change and when an exception should end, we also train a lightweight linear
sub-task probe on frozen VLA features. Its predictions supplement a grasp state machine driven by execution feedback;
the state machine determines attachment and detachment. An optimization layer uses the resulting constraints to
refine the portion of each action chunk that will be executed. The VLA remains frozen after embodiment fine-tuning:
\emph{no additional VLA training for collision avoidance is required}. The probe is trained separately, and the current
implementation reads attention through an additional pass of the same checkpoint.

The contribution is the connection from internally derived target geometry to phase- and link-specific constraints.
We develop it through:
\begin{itemize}
\item An attention-to-geometry procedure that extracts an observed 3D target from a frozen VLA and calibrated depth,
without an external segmentation model or a separately supplied target name.
\item A contact representation that uses this target geometry for link-specific distance queries, manages its lifetime
with sub-task and grasp-state information, and transfers grasped geometry to the robot's collision queries.
\item \tbd{A controlled evaluation: insert verified task success, forbidden-contact and safe-success results, then state
what target grounding and link-specific handling add under the same refinement controller.}
\end{itemize}
The evaluation separates target identity, geometric extent, and contact handling so that improved execution can be
attributed to the information supplied to the controller rather than to a change of controller alone.
